\scanpage{1}
% not manuscript: printed Montpellier archival cover for cote 140-3.

\scanpage{2}
\unsure{Pages} 291--650 $+$ table des matières

\uline{NB} Les pages 291,292 figurent deux fois
(erreur de numérotation)

\margin{1}

\scanpage{3}
\uline{$2^{\text{ème}}$ partie}\quad
\fbox{Variations sur l'équation des lacets}

\uline{38}\quad Détermination de
$\Pi'_{0,3}\xrightarrow{\sim}\unsure{\mathcal C|\mathcal G}_{0,3}^{+}$.
\hfill 291

\uline{39}\quad Propriétés de congruence \unsure{de|des} l'\unsure{opération|opérations}
de $\unsure{\Gamma|\Pi}_{\mathbb Q}$ sur $\widehat{\Pi}_{0,3}$,
\unsure{Lien} : l'hom
$N_{\unsure{1,1}}\longrightarrow \mathrm{Gl}(\unsure{1|2},\widehat{\mathbb Z})$
\hfill \unsure{299|297}

40\quad Paraphrase \unsure{topologico-isotopique} de la
théorie de Legendre
\hfill $307'$

41\quad Conditions de commutation \unsure{de|des}
$\widehat{\Pi}'_{0,3}\xrightarrow{\sim}\widehat{\unsure{\mathcal C|\mathcal G}}_{0,3}^{+}$
aux op. ext. de $\unsure{\Gamma|\Pi}_{\mathbb Q}$.
\hfill \deleted{\illegible{1}}

Calculs des \unsure{autom.} ext. des $\widehat{\Pi}_{0,3}$ et
\unsure{des} $\widehat{\unsure{\mathcal C|\mathcal G}}_{0,3}^{+}$.
\hfill 315

\uline{42}\quad Récapitulation sur le groupe
$\unsure{\mathcal C|\mathcal G}_{0,3}$
$\unsure{\simeq\ldots}$ et sur les autom. à lacets des
$\widehat{\unsure{\mathcal C|\mathcal G}}_{0,3}^{+}$ et des
$\widehat{\Pi}_{0,3}$ :
\hfill 331

\[
\left\{
\begin{array}{llr}
\mathrm{I}&\unsure{\mathcal T|\mathcal C}_{0,3}&331\\
\mathrm{II}&\Pi_{0,3}\text{ et }\unsure{\mathcal C|\mathcal G}_{0,3}&332\\
\mathrm{III}&\text{relations avec }\Pi'_{0,3}\quad
\deleted{\text{et }\mathrm{Sl}(2,\mathbb Z)'}&\unsure{334}\\
\mathrm{IV}&\deleted{\text{\unsure{homomorphismes de}}\ \unsure{\mathcal C|\mathcal G}_{0,3}}\quad
\inserted{\text{relations avec }\mathrm{Sl}(2,\mathbb Z)'}&\unsure{335'}\\
\mathrm{V}&\text{les \unsure{acteurs} de }\unsure{\mathcal C|\mathcal G}_{0,3}&336'
\end{array}
\right.
\]

VI\quad \unsure{Autom.} à lacets des $\widehat{\Pi}_{0,3}$ et
$\widehat{\unsure{\mathcal C|\mathcal G}}_{0,3}^{+}$
\hfill $337'$

VII\quad Calculs dans $\mathrm{Sl}(2,\widehat{\mathbb Z})'$ :
cocycles \unsure{transp}
\hfill 342

VIII\quad Calculs dans $\mathrm{Sl}(2,\widehat{\mathbb Z})'$ :
cocycles \unsure{non} $\unsure{\sigma|\overline{\sigma}}_{\unsure{\infty|0}}$
\hfill 346

\margin{\uline{NB} les calculs de ces \unsure{sections} \uline{VI},
\uline{VII}, \uline{VIII} en partie \unsure{remplacés} par ceux
du \unsure{\S46}, \unsure{avec} des notations plus simples\ldots}

43\quad Digression sur les $1$-cocycles des groupes \unsure{cycl.}
\hfill 350

44\quad Retour sur notations\inserted{(les $\unsure{\rho|\beta|g},\unsure{h},f,\alpha,\beta\ldots$)}.
Mise en équations dans $\mathrm{Sl}(2,\widehat{\mathbb Z})$.
\hfill 356

45\quad L'hom.
$\widetilde{M}_{0,3}\longrightarrow \mathrm{Gl}(2,\widehat{\mathbb Z})'$,
et construction de l'extension $\unsure{\widetilde{N}_{0,3}}$
de $\widetilde{\Gamma}$ par $\widehat{\mathrm{Sl}(2,\mathbb Z)}$.
\hfill \unsure{384}

\uline{46}\quad Notations définitives (?) -- récapitulation
des calculs \unsure{sur} les \unsure{autom.} à lacets
\unsure{des} $\widehat{\Pi}_{0,3}$ et de
$\widehat{\unsure{\mathcal C|\mathcal G}}_{0,3}^{+}$
\hfill 390

\margin{2}

\scanpage{4}

I Les groupes \(M_{0,3}^{\sim}\) et \unsure{\(N_{0,3}^{!\sim}\)} \hfill 390

\uline{II} Calcul des éléments de \(N_{0,3}^{!\sim}\) \hfill 392

1) Via \(f,g,h\)

2) Via \(\alpha_0,\beta_0,f\) \hfill 392

3) Via \(\alpha,\beta,f\) \hfill 393

4) Via \unsure{\(\zeta,\eta,f\)} \hfill \(393'\)

5) Comparaison \unsure{entre} \illegible{1}, \unsure{relations} des \hfill 394
cocycles \hfill 395

\uline{III} Calculs dans \unsure{\(\mathrm{Sl}(2,\widehat{\mathbb Z})'\)} et dans \(\mathrm{Gl}(2,\widehat{\mathbb Z})\) \hfill 396

\uline{IV} Les trois cas fondamentaux, \deleted{et les}
\unsure{\((\alpha_0=1,\beta_0=1,f=1)\)}, et les \unsure{relèvements}
\(K_{1/2},K_2,K_{-1}\). Les sous-groupes \(M_{0,3}^{!\sim}(q)\),
\(q\in\{\)\deleted{\unsure{\(0,1,\infty\)}}\inserted{\unsure{\(-j,-\overline j\)|\(-j,-j^2\)}}\(,2,-1,1/2\}\)
\hfill \deleted{\unsure{402}} 402

\uline{V} Récapitulation en termes de \(M_{0,3}^{!\sim}[0]\)
\hfill \deleted{\unsure{404}}

Les sous-groupes \(M_{0,3}^{!\sim}(q)\), \(q\notin\{0,1,\infty\}\) \hfill 407

\margin{calcul \unsure{en §17}}
\uline{VI} Retour sur l'introduction de \(\alpha\) et \(\beta\) en termes de \(\zeta,\eta\). \hfill 414

47. Retour sur l'homomorphisme
\unsure{\(\widehat\Pi_{0,3}\xrightarrow{\widehat\varphi}\widehat{\mathcal C}_{0,3}^{+}\)},
les groupes \unsure{\(M_{0,3}^{\approx},N_{0,3}^{\approx},M_{0,3}^{!\approx}(0)\)|\(M_{0,3}^{\sim},N_{0,3}^{\sim},M_{0,3}^{!\sim}(0)\)}
\hfill \deleted{\unsure{415}} 415

\uline{48} Construction d'un schéma en \unsure{groupes} \hfill 420

I) Étude du morphisme \(\alpha\mapsto\alpha\rho(\alpha)^{-1}=\)\unsure{\([\alpha,\rho]\)} \hfill 420

II) L'équation \(\zeta=\eta\varepsilon_1(\mu-1)\), et le groupe
\unsure{\(B'_0\simeq B_{\rho,\sigma}\)} \hfill \(434'\)

\uline{III}) L'équation \(\alpha\rho(\alpha)^{-1}=\beta\sigma(\beta)^{-1}\varepsilon_1(\mu-1)\),
et le groupe \unsure{\(\mathcal G_{\rho,\sigma}^{c}\)} de \illegible{1} \hfill \(443'\)

(IV) Un épinglage canonique associé à \(\rho,\sigma\) \hfill 445

(V) L'équation \unsure{\(\alpha\rho(\alpha)^{-1}=\varepsilon(\beta\sigma^{-1}(\beta)^{-1})\varepsilon_1(\mu-1)\)|\(\alpha\rho(\alpha)^{-1}=\varepsilon(\beta\sigma'(\beta)^{-1})\varepsilon_1(\mu-1)\)}
\hfill \deleted{\unsure{454}}

\unsure{\(\varepsilon\)} Son étude en
\unsure{\(\widetilde\alpha\rho(\widetilde\alpha)^{-1}=\widetilde\beta\sigma(\widetilde\beta)^{-1}\varepsilon_1(\mu-1)\)} --
\inserted{dans \unsure{\((\mathbb Q_p(i))\)}}
le groupe \unsure{\(\mathcal G_{\rho,\sigma}^{*}\)}, la condition \(\mathrm{Tr}\,\sigma=0\) \hfill 459

(VI) Caractérisation \inserted{modulaire} du triplet \unsure{\((\varepsilon_3,\sigma,\rho)\)|\((\varepsilon,\sigma,\rho)\)} \hfill 462

\uline{VII} Récapitulation sur les épinglages. \hfill \unsure{476}

\uline{VIII} Récapitulation sur l'équation -- le groupe
\unsure{\(S\mathcal G_{\rho,\sigma}^{*}\)} et ses variantes \hfill \unsure{483}

\uline{IX} Les groupes \(M_{\rho,\sigma}\) et \(S_0M_{\rho,\sigma}\), et les homomorphismes
\unsure{\(S_0M_{0,3}^{!\sim}\longrightarrow SM_{\rho,\sigma}(\widehat{\mathbb Z})/\mathbb Z\)|\(S_0M_{0,3}^{!\sim}\longrightarrow SM_{\rho,\sigma}(\widehat{\mathbb Z}/\mathbb Z)\)},
\unsure{\(M_{0,3}^{!\sim}\longrightarrow M_{\rho,\sigma}(\widehat{\mathbb Z})/\mathbb Z\)} \hfill \(490'\)

\uline{X} Structure de \(M_{\rho,\sigma}\), et l'homomorphisme canonique
\(M_{\rho,\sigma}\longrightarrow\) \unsure{\(G\)}\deleted{\illegible{1}}\(\simeq\mathrm{Gl}(2)\) \hfill 499

\uline{XI} Récapitulation sur \(M_{\rho,\sigma}\).
\deleted{les deux systèmes de paramètres \unsure{\((\underline U,\underline a,\underline b)\)} et
\((\alpha,\beta,f)\), sur \unsure{\(M_{\rho,\sigma}[0]\)}}
\hfill \deleted{\illegible{1}}

\deleted{\uline{XII} Les sous-groupes \(M_{\rho,\sigma}[0]\), \unsure{\(M_{\rho,\sigma}(1)\)}, \illegible{1}}

\(1^{\circ})\) Données préliminaires \hfill 508

\(2^{\circ})\) \deleted{\unsure{\(M_{\rho,\sigma}(\mathbb Z)\)}}
\(M_{\rho,\sigma}[0],M_{\rho,\sigma}(\infty)\),
\unsure{\(\mathcal T_{\rho,\sigma},\mathcal T_{\rho,\sigma}^{0}\)},
\unsure{\(\mathcal M_{\rho,\sigma},\mathcal C_{\rho,\sigma}^{+},\mathcal T_{\rho,\sigma}^{+0}\)}
\hfill \deleted{\unsure{524}} 509

\margin{\unsure{repris au} \(6^{\circ})\)}
[\(3^{\circ})\) \(M_{\rho,\sigma}^{0},M_{\rho,\sigma}^{+},\ldots\)
\unsure{calculs} de pts entiers] \hfill 511

\(4^{\circ})\) Les hom \(\Theta:M_{\rho,\sigma}\longrightarrow G\),
\(\overline\Theta:M_{\rho,\sigma}\longrightarrow PG\) \hfill \(512'\)

\(5^{\circ})\) \(R_{\rho,\sigma}^{0},SR_{\rho,\sigma}^{0}\)] \hfill \unsure{\(513'\)}

\uline{\(6^{\circ})\)} \unsure{Groupes}, \illegible{3} et \(R_{\rho,\sigma}\).
\unsure{Centres} de \(M_{\rho,\sigma},M_{\rho,\sigma}^{0}\).
Le ss-groupe \(\Sigma\).
\(\Sigma\simeq\mathbb Z/2\mathbb Z*\mathbb Z/2\mathbb Z*\mathbb Z/2\mathbb Z\).
\hfill \deleted{\unsure{514|517}} 517

\uline{\(7^{\circ})\)} Dévissage de \(R_{\rho,\sigma}\) et de \(M_{\rho,\sigma}\) \hfill \(526'\)

\(8^{\circ})\) Les paramètres
\unsure{\(\underline U,\underline a,\underline b,\mu,\varepsilon,\varepsilon'\)} et
\(\underline u,f,\alpha,\beta,\alpha_0,\beta_0\)
\unsure{\((\zeta,\eta,g,h)\)} \hfill 529

\(9^{\circ})\) Les sous groupes \(M_{\rho,\sigma}(0),M_{\rho,\sigma}[0]\), et le groupe
\(SM_{\rho,\sigma}\), l'isom. \unsure{\(SM_{\rho,\sigma}\simeq S\Sigma_{\rho,\sigma}^{*}\)} \hfill 534

\(10^{\circ})\) Les sous-groupes \inserted{section}
\unsure{\(M_{\rho,\sigma}(\Phi),M_{\rho,\sigma}(P)\)} \hfill \unsure{\(540'\)|\(550'\)}

\uuline{XII} \uline{Rétrospective.} \hfill 542

\margin{3}

\scanpage{5}

\uline{49}. Les groupes \(M_{\rho,\sigma}\) généraux \hfill 550

[I Préliminaires heuristiques]

\uline{II} Les groupes \unsure{\(M_{\rho,\sigma}^{*}\)} \unsure{dits|et} ``sans signe'' -- nouvelles
variations sur la relation des lacets \hfill \unsure{\(553'\)}

\uline{III} Passage de \(M_{0,3}^{\sim}\) à \(N_{1,1}^{\sim}\) -- et de
\unsure{\(\mathrm{Gl}(2,\mathbb Z)^{\wedge}/\pm1\)|\(\mathrm{Gl}(2,\mathbb Z)^{\sim}/\pm1\)} à \(\mathrm{Gl}(2,\mathbb Z)^{\wedge}\) \hfill 565

[\uline{IV} Où on se bat avec ``les signes'' \(\varepsilon_2,\varepsilon_3\)
dans la relation des lacets -- calculs -- brouillons\ldots] \hfill 568 \deleted{\unsure{577}}

\deleted{\illegible{1}}\textcircled{\uline{V}} Les groupes \(M_{\rho,\sigma}\) généraux
(contexte \uline{profini}) \hfill 577

\(1^{\circ})\) Données \hfill 577

\(2^{\circ})\) \(Z_{\rho,\sigma}\), \unsure{\(\sigma_{\rho,\sigma}\)|\(\delta_{\rho,\sigma}\)} \hfill \(577'\)

\(3^{\circ})\) \(S_0G_{\rho,\sigma},G_{\rho,\sigma}\) \hfill 579

\(4^{\circ})\) \(Z_\rho,Z_\sigma,N_\rho,N_\sigma\),
\unsure{\(N_\rho^{*},N_\sigma^{*},SN_\rho^{*},SN_\sigma^{*}\)} \hfill - - - 580

\(5^{\circ})\) \(S_0M_{\rho,\sigma},M_{\rho,\sigma},S_0\Gamma_{\rho,\sigma},\Gamma_{\rho,\sigma}\)
\hfill - - - - - - 585

\(6^{\circ})\) \unsure{\(\Sigma_{\rho,\sigma},S\mathcal G_{\rho,\sigma},\Sigma_{\rho,\sigma}^{\mathrm{eff}},S\mathcal G_{\rho,\sigma}^{\mathrm{eff}}\)}
\hfill 586

\(7^{\circ})\) Fonctorialités. \hfill 588

\uline{\(8^{\circ})\)} Le cas universel, \unsure{\(\mathcal C^{\wedge}\)}\(=\mathrm{Gl}(2,\mathbb Z)^{\wedge}\).
Le \deleted{\illegible{1}} sous-groupe \(\Pi_0\). Étude de l'extension
\(S_0\Gamma_{\rho,\sigma}^{!}\) de \(Z_{\rho,\sigma}^{!}=M^{!}[0]\) par
\(L_\rho*L_\sigma\), et \unsure{\(S_0^{\Pi}\Gamma_{\rho,\sigma}^{!}\)} de
\(Z_{\rho,\sigma}^{!}\) par \(\mu=\{\pm1\}\).
\hfill \deleted{\unsure{590}} 590

\margin{\unsure{rectificat.} de la ``relation \unsure{de conjugaison}'' \illegible{1} \unsure{liée} \deleted{\illegible{1}}}

\(9^{\circ})\) Récapitulation sur le cas universel.
Structure des \unsure{\(N_\sigma^{\circ},N_\rho^{*}\)}. L'hom can.
\unsure{\(M\xrightarrow{\sim}M_{\rho,\sigma}^{\unsure{\rightleftarrows|\leftrightarrows}}/H_\rho\hookrightarrow M_{\rho,\sigma}/H_\rho\)}.
\hfill \deleted{\illegible{1}} 614

\deleted{\illegible{1}}\textcircled{\uline{VI}} Les groupes \(M_{\rho,\sigma}\) généraux:
\unsure{récapitulation}, et raffinement de la construction précédente. \hfill 628

\(1^{\circ})\) Données (628) \quad
\(2^{\circ})\) \(Z_{\rho,\sigma}\), \unsure{\(\sigma_{\rho,\sigma}\)|\(\delta_{\rho,\sigma}\)} (629)

\(3^{\circ})\) \unsure{\(G_{\rho,\sigma},G_{\rho,\sigma}^{\unsure{\rightleftarrows|\leftrightarrows}}\)}\ldots\ (\unsure{632|631}) \quad
\(4^{\circ})\) \unsure{\(N_\rho,N_\sigma,N_\rho^{\unsure{\rightleftarrows|\leftrightarrows}},N_\sigma^{\unsure{\rightleftarrows|\leftrightarrows}}\)}\ldots\ (\unsure{635})

\(5^{\circ})\) \unsure{\(M_{\rho,\sigma},M_{\rho,\sigma}^{\unsure{\rightleftarrows|\leftrightarrows}}\)}\ldots\ (639) \quad
\(6^{\circ})\) \unsure{Fonctorialités|Fonctorialité} (\(640'\))

\margin{\(7^{\circ})\) \unsure{\(M_{\rho,\sigma}[\infty],M_{\rho,\sigma}[\rho],M_{\rho,\sigma}[\sigma]\)|\(M_{\rho,\sigma}[0],M_{\rho,\sigma}[\rho],M_{\rho,\sigma}[\sigma]\)} (642)

\(8^{\circ})\) Retour sur le cas universel: \unsure{\(M^{\unsure{\rightleftarrows|\leftrightarrows}}\)} et l'hom
\unsure{\(M^{\unsure{\rightleftarrows|\leftrightarrows}}\longrightarrow M_{\rho,\sigma}^{\unsure{\rightleftarrows|\leftrightarrows}}/H_\rho\)} (\(643'\))

\(9^{\circ})\) Digression de \illegible{1} \unsure{formalisme} des groupes 648

\((10^{\circ})\) \unsure{Lien|Liens} avec §48}

\margin{4}

\scanpage{6}
\mspage{\unsure{291bis}}

\uline{38}\quad \unsure{Détermination} explicite de l'hom
$\Pi'_{0,3}\xrightarrow{\sim}\unsure{\mathcal C|\mathcal G}_{0,3}^{+}$
(et de $\unsure{\Pi_{0,3}^{\prime\sim}}$
(gr. cartographique triangulé non orienté)
$\xrightarrow{\sim}\unsure{\mathcal C|\mathcal G}_{0,3}$).

J'ai envie de revenir sur la prop. du \S\ précédent
(p.\ 262) avec surface \inserted{\unsure{connexe}} $X$, groupe discret $G$
opérant proprement, les stabilisateurs \inserted{$G_x$} des $x\in X$
\deleted{\unsure{étant}} conservant l'orientation et opérant fidèlement
au voisinage de $x$ \ldots\ en prenant $Y=X/G$, \unsure{muni|munie} de la
``signature'' $\underline d=d(X/Y)$ définie par la \unsure{ramification}
--- d'où un \uline{\unsure{Iso|iso}}
\[
(1)\qquad
\Pi=\Pi_1((X,G),x)\xrightarrow{\sim}\Pi_1((Y,\underline d),y)=\Pi'
\]
($x\in X$, $y$ image de $x$ dans $Y$)
que je voudrais \unsure{étudier} de plus près.

On a dans $\Pi_1((Y,\underline d),y)$, pour tout
$b\in\unsure{\operatorname{Supp}}\underline d=Y^!$,
des ``sous-groupes \unsure{lacet-tronqué}'', plus précisément
une \unsure{classe} de conjugaison \unsure{de|des} \unsure{ss-groupes}
quotients de $\mathbb Z/d_b\mathbb Z$, et plus précisément encore,
un hom extérieur canonique
\[
(2)\qquad
\unsure{T_b|T_y}\otimes_{\mathbb Z}\mathbb Z/d_b\mathbb Z
\longrightarrow\Pi'
\]
déduit par passage au quotient de l'hom ext
\[
(3)\qquad
\unsure{T_b|T_y}\xrightarrow{\unsure{k_y|k_b}}
\Pi_1(Y\setminus Y^!)
\qquad(\longrightarrow\Pi_1((Y,\underline d)))
\]
(où $T_b$ est le $\mathbb Z$-module libre de rang $1$ des
``entiers tordus'' en $b$ sur $Y$). D'autre part,
si $a\in X$ est au dessus de $b$, le stabilisateur $G_a$ est
\deleted{\unsure{donc}} un groupe cyclique d'ordre $d_a=d_b$
fini --- On a \unsure{$=$} un isom. can. évident
\[
(4)\qquad
G_a\simeq T_a\otimes_{\mathbb Z}\mathbb Z/d_a\mathbb Z
\simeq T_b\otimes\mathbb Z/d_b\mathbb Z.
\]
\margin{(NB $T_a\simeq T_b$ \unsure{canoniq\textsuperscript{t}})}
\margin{5}

\scanpage{7}

D'autre part, à $a$ est \unsure{associé|associée} une classe de
conjugaison de \unsure{scindages} de l'extension
\[
(5)\qquad
1\longrightarrow\Pi_1(X,x)\longrightarrow\Pi_1((X,G),x)
\longrightarrow G\longrightarrow1,
\]
(cf \S\ 15)
\[
(6)\qquad
k_a:G_a\longrightarrow\Pi_1((X,G),x)
\]
ces classes \inserted{de $\Pi_1(X,\unsure{x|e})$-conj.}
de \unsure{scindages} (si $X\ne S^2$) correspondent
biunivoquement aux pts fixes de $\unsure{G_0|G_a}$ opérant sur $X$.
Si on prend l'\uline{hom} \uline{extérieur} \unsure{défini|définis}
par (6) (donc en \unsure{conjuguant|conjugant} non seulement par les
$\in\Pi_1(X)$, mais par ceux de $\Pi_1(X,G)$ tout entier)
on trouve une corr. $1$-$1$ entre \inserted{l'ens.}
des hom ext. (6) qui relèvent l'inclusion
\[
\unsure{G_0|G_a}\hookrightarrow G
\]
(\unsure{considérée} comme hom extérieur), et l'ens des orbites de $G$
dans $X^!$, i.e. entre les pts de $Y^!$. En \unsure{résumé}

\uline{Proposition}\quad On a une bijection entre $Y^!$
et l'ens des hom ext. ``maximaux''
\inserted{$\varphi=\varphi_b:\Gamma\longrightarrow\unsure{\Pi|\pi}$}
de \unsure{ss-groupes} $\Gamma\subset G$, \deleted{\unsure{finis}}
$\Gamma\ne1$, dans $\Pi_1(X,G)$ qui relèvent l'inclusion
\inserted{$\Gamma\subset G$} (les dits $\Gamma$ sont nécessairement
cycliques, les $\Gamma$
\margin{6}

\scanpage{8}
\mspage{292bis}

en question \unsure{correspondant|correspondants} à un $b\in Y^!$
étant les stabilisateurs $G_a$ dans $G$ des
\deleted{\unsure{$a\in(X^!)_a$|$a\in(X^!)_b$}}
$a\in X_b=X_b^!$. De plus, le composé des hom.
\[
(7)\qquad
T_b\otimes\mathbb Z/\unsure{n|m}\mathbb Z
\overset{(4)}{\simeq}G_a
\xrightarrow{\varphi_a=\varphi_b}\Pi_1(X,G)
\overset{(1)}{\xrightarrow{\sim}}\Pi_1((Y,\underline d))
\]
n'est autre (on s'en serait douté !) que \unsure{``}l'hom
``\unsure{lacets tronqué}'' (2).

Supposons maintenant que $X$ soit de la forme
$\widehat X\setminus S$, $\widehat X$ surface compacte,
$S$ partie finie de $\widehat X$ \deleted{\unsure{donc}}
\[
(8)\qquad
\left\{\begin{aligned}
X&=\widehat X\setminus S\quad\text{donc}\quad
Y=\widehat Y\setminus T\quad\unsure{\ne\varnothing}\quad\text{\unsure{avec}}\\
\widehat Y&=\widehat X/G
\end{aligned}\right.
\]

Considérons\unsure{,} alors, pour tout $\unsure{a|b}\in S$, l'hom extérieur
\[
(9)\qquad
k_a:T_a\longrightarrow\Pi_1(X)\hookrightarrow\Pi_1(X,G)
\]
que je vais comparer au composé
\[
(10)\qquad
\deleted{\unsure{k_b:}}\quad T_b\xrightarrow{k_b}\Pi_1(Y)
\longrightarrow\Pi_1(Y,\underline d)
\]
(pour $b=\unsure{\widehat f|\widehat j}(a)$), moyennant l'isom. can.
\[
\deleted{\unsure{(12)}}\ (11)\qquad T_a\simeq T_b
\]
Si $d_a$ désigne encore l'indice de ramification en $a$, on aura
un diagramme commutatif de compatibilité:
\margin{7}

\scanpage{9}

\[
(12)\qquad
\diagram{\begin{tikzcd}[ampersand replacement=\&,column sep=large,row sep=large]
T_a \arrow[r,"k_a"] \arrow[d,"\wr\ \mathrm{can}"'] \& \Pi_1(X) \arrow[r] \& \Pi_1(X,G) \arrow[dd,"\wr"] \\
T_b \arrow[d,"d_a\cdot\mathrm{id}"'] \& \& \\
T_b \arrow[rr,"k_b"] \& \& \Pi_1(Y,\underline{d})
\end{tikzcd}}
\]

Nous allons appliquer ceci à la \unsure{situation} de
$X=U_{0,3}^{\mathrm{an}}$, $G=\unsure{\mathfrak{S}|\mathcal{S}}_3$,
$\widehat{X}=\mathbb{P}^{1\mathrm{an}}$, $\widehat{Y}\simeq\mathbb{P}^{1\mathrm{an}}$,
\unsure{avec} l'application de \unsure{passage} au quotient
\unsure{donné|donnée} par (p. 286' (58))
\[
(13)\qquad
f(z)=\frac{3^3}{2^2}\frac{z^2(z-1)^2}{(z^2-z+1)^3}
\]
\[
f:\quad
\left\{\begin{array}{l}0\\1\\\infty\end{array}\right.\longrightarrow 0
\qquad
\left\{\begin{array}{l}2\\-1\\1/2\end{array}\right.\longrightarrow 1
\qquad
\left\{\begin{array}{l}-\overline{j}\\-j\end{array}\right.\longrightarrow\infty
\]

On voit donc, avec les notations de \hspace{3em}, que les générateurs
$l'_0,l'_1,l'_\infty$ de $\Pi'_{0,3}$
(\unsure{satisfaisant}
\[
(14)\qquad l'_\infty l'_1 l'_0=1
\qquad {l'_1}^{2}={l'_\infty}^{3}=\unsure{\cdot}1\qquad )
\]
\unsure{correspondent}, respectivement, par l'isomorphisme
\[
(15)\quad
\begin{array}{ccl}
\unsure{\mathcal{C}|\mathcal{G}}_{0,3}^{+}
=\Pi_1(U_{0,3}^{\mathrm{an}},\unsure{\mathfrak{S}|\mathcal{S}}_3)
&\xrightarrow{\sim}&\Pi'_{0,3}=\mathcal{T}_{1,1}^{+\prime}\\
\begin{array}{c}
\text{groupe cartographique}\\
\text{triangulé \unsure{pondéré} or.}\\
\text{\uline{mixte}}
\end{array}
&&\begin{array}{c}
\text{groupe cartographique \unsure{triangulé|triangul-}}\\
\text{orienté}\qquad \simeq\mathrm{Sl}(2,\mathbb{Z})'
\end{array}
\end{array}
\]
aux éléments notés
\deleted{$\unsure{\rho,\varepsilon_0,\varepsilon_1}$ \illegible{2}}
\[
(16)\qquad
\left\{\begin{array}{lll}
\varepsilon_0,\varepsilon_1,\varepsilon_\infty&\longmapsto&l'_0\\
\text{(tels que $\varepsilon_i^2=l_i$)}&&\\
\sigma_0,\sigma_1,\sigma_\infty&\longmapsto&l'_1\\
\rho&\longmapsto&l'_\infty
\end{array}\right|\quad
\begin{array}{l}\mathrm{mod}\\\text{\unsure{conjugaison}}\end{array}
\]
\margin{ou $\rho^{-1}$ ?}
\margin{8}

\scanpage{10}
\mspage{293}

\deleted{\unsure{L'isom}}

Ceci montre que l'on peut trouver
\inserted{\unsure{dans|de} $\Pi=\unsure{\mathcal{C}|\mathcal{G}}_{0,3}^{+}$}
des conjugués $\varepsilon'_0$ de $\varepsilon_0$, $\sigma'_0$ de $\sigma_0$,
tels qu'on ait
\[
(17)\qquad
\rho\sigma'_0\varepsilon'_0=1,\quad
\text{en plus de}\quad \rho^3={\sigma'_0}^{2}=1
\]
et de telle façon que ces relations \unsure{fournissent}
une \uline{présentation} de $\Pi$ (donc en particulier,
que $\varepsilon'_0,\sigma'_0,\rho$ \unsure{engendrent} \unsure{$\}$} $\Pi$).
S'inspirant d'ailleurs de l'isomorphisme \deleted{plus} qui
\unsure{prolonge} (15) et s'obtient à \unsure{l'avenant}
\[
(16)\qquad
\unsure{\mathcal{C}|\mathcal{G}}_{0,3}
\simeq\Pi_1\left(U_{0,3}^{\mathrm{an}},
\underbrace{\unsure{\mathfrak{S}|\mathcal{S}}_3\times\mathbb{Z}_2}_{\mathcal{T}_{0,3}}\right)
\longrightarrow
\begin{array}{c}
\mathcal{T}'_{1,1}\\
\text{groupe cartographique}\\
\text{triangulé non}\\
\text{orienté}\\
\simeq\mathrm{Gl}(2,\mathbb{Z})
\end{array}
\]
on devrait trouver dans le premier \illegible{1}
trois \unsure{générateurs|generateurs} $\tau'_0,\tau'_1,\tau'_\infty$
(\unsure{correspondants} aux gén. \unsure{analogues}
$\sigma_0,\sigma_1,\sigma_\infty$ du groupe cartographique
\unsure{correspondant}) \unsure{satisfaisant|satisfaisants}
\[
(17)\qquad
{\tau'_0}^{2}={\tau'_1}^{2}={\tau'_\infty}^{2}=1,\quad
\left(\underbrace{\tau'_{\unsure{0|\infty}}\tau'_\infty}_{\begin{array}{c}\text{\unsure{corr.} à}\\l'_1\end{array}}\right)^2
=\left(\underbrace{\tau'_1\tau'_0}_{\begin{array}{c}\text{corr. à}\\l'_\infty\end{array}}\right)^3=1
\]
de telle façon que ce soient là \deleted{\unsure{des}} des
relations de présentation de $\unsure{\mathcal{C}|\mathcal{G}}_{0,3}$.
Il faudrait que
\deleted{$\tau'_\infty\tau'_0$ soit conjugué de $\sigma_0$}
\inserted{l'on ait}
\[
(18)\qquad
\begin{array}{rclcrclcrcl}
\tau'_\infty\tau'_1&=&\varepsilon'_0;&\quad&
\tau'_0\tau'_\infty&=&\sigma'_0,&\quad&
\tau'_1\tau'_0&=&\rho\\
&&\updownarrow&&&&\updownarrow&&&&\updownarrow\\
&&l'_0&&&&l'_1&&&&l'_\infty
\end{array}
\]
\margin{9}
